\documentclass[11pt,fontset=fandol]{ctexart}
\usepackage[a4paper,margin=23mm]{geometry}
\usepackage{amsmath,amssymb,mathtools,bm,mathrsfs,esint}
\usepackage{graphicx,longtable,booktabs,array,calc}
\usepackage[normalem]{ulem}
\usepackage{xcolor,hyperref}
\hypersetup{hidelinks}
\setcounter{secnumdepth}{-1}
\newcounter{none}
\providecommand{\tightlist}{\setlength{\itemsep}{0pt}\setlength{\parskip}{0pt}}
\providecommand{\passthrough}[1]{#1}
\usepackage{amssymb,mathrsfs,esint,bm,newunicodechar}
\allowdisplaybreaks
\setlength{\emergencystretch}{3em}
\providecommand{\llangle}{\langle\!\langle}
\providecommand{\rrangle}{\rangle\!\rangle}
\newunicodechar{∞}{\ensuremath{\infty}}
\newunicodechar{∈}{\ensuremath{\in}}
\newunicodechar{∑}{\ensuremath{\sum}}
\newunicodechar{∫}{\ensuremath{\int}}
\newunicodechar{≤}{\ensuremath{\leq}}
\newunicodechar{≥}{\ensuremath{\geq}}
\newunicodechar{≠}{\ensuremath{\neq}}
\newunicodechar{→}{\ensuremath{\to}}
\newunicodechar{⇒}{\ensuremath{\Rightarrow}}
\newunicodechar{α}{\ensuremath{\alpha}}
\newunicodechar{β}{\ensuremath{\beta}}
\newunicodechar{γ}{\ensuremath{\gamma}}
\newunicodechar{λ}{\ensuremath{\lambda}}
\newunicodechar{φ}{\ensuremath{\varphi}}
\newunicodechar{ε}{\ensuremath{\varepsilon}}
\newunicodechar{σ}{\ensuremath{\sigma}}
\newunicodechar{θ}{\ensuremath{\theta}}
\newunicodechar{δ}{\ensuremath{\delta}}
\newunicodechar{η}{\ensuremath{\eta}}
\newunicodechar{π}{\ensuremath{\pi}}
\newunicodechar{ξ}{\ensuremath{\xi}}
\newunicodechar{Σ}{\ensuremath{\Sigma}}
\newunicodechar{Δ}{\ensuremath{\Delta}}
\newunicodechar{∂}{\ensuremath{\partial}}
\newunicodechar{∀}{\ensuremath{\forall}}
\newunicodechar{∃}{\ensuremath{\exists}}
\newunicodechar{√}{\ensuremath{\surd}}
\newunicodechar{−}{\ensuremath{-}}
\newunicodechar{μ}{\ensuremath{\mu}}
\newunicodechar{ρ}{\ensuremath{\rho}}
\newunicodechar{τ}{\ensuremath{\tau}}
\newunicodechar{ω}{\ensuremath{\omega}}
\newunicodechar{Ω}{\ensuremath{\Omega}}
\newunicodechar{χ}{\ensuremath{\chi}}
\newunicodechar{ψ}{\ensuremath{\psi}}
\newunicodechar{Γ}{\ensuremath{\Gamma}}
\newunicodechar{Φ}{\ensuremath{\Phi}}
\newunicodechar{Λ}{\ensuremath{\Lambda}}
\newunicodechar{ℝ}{\ensuremath{\mathbb{R}}}
\newunicodechar{ℤ}{\ensuremath{\mathbb{Z}}}
\newunicodechar{ℕ}{\ensuremath{\mathbb{N}}}
\newunicodechar{△}{\ensuremath{\triangle}}
\newunicodechar{⊗}{\ensuremath{\otimes}}
\begin{document}
\section{2025年高等代数真题}\label{ux5e74ux9ad8ux7b49ux4ee3ux6570ux771fux9898}

来源：\texttt{HDU/601数学分析/2025/2025年杭州电子科技大学数分和高代真题(双科合卷).pdf}。原
PDF 页序：3, 4。

\section{杭州电子科技大学2025年高等代数真题}\label{ux676dux5ddeux7535ux5b50ux79d1ux6280ux5927ux5b662025ux5e74ux9ad8ux7b49ux4ee3ux6570ux771fux9898}

\subsection{1.（15分）}\label{ux5206}

证明：\(1+\sqrt3+\sqrt[3]9\) 为无理数。

\subsection{2.}\label{section}

设 \(\alpha,\beta,\gamma\)
为三维列向量，\(A_1=(\alpha,\beta,\gamma)\)，\(|A_1|=-2\)，且
\(A_2=(5\alpha,-3\beta,2\gamma)\)，求

\[
\begin{vmatrix}
0&A_1\\
A_2&0
\end{vmatrix}.
\]

\subsection{3.}\label{section-1}

证明：\(A_{m\times n}X=0\) 和 \(B_{s\times n}X=0\) 同解的充要条件为

\[
r(A)=r(B)=r\begin{pmatrix}A\\B\end{pmatrix}.
\]

\subsection{4.}\label{section-2}

已知 \(A\) 为正定矩阵，证明：\(A+A^{-1}-2E\) 为半正定矩阵。

\subsection{5.}\label{section-3}

若

\[
A=\begin{pmatrix}
-10&5&-1\\
-12&5&1\\
-12&1&9
\end{pmatrix}.
\]

（1）是否有 \(A=B^3\)？

（2）将（1）进行推广，并证明。

\subsection{6.}\label{section-4}

设二次型

\[
f(x_1,x_2,x_3)=3x_1^2+6x_2^2+3x_3^2-4x_1x_2-8x_1x_3-4x_2x_3,
\]

在条件 \(x_1^2+x_2^2+x_3^2=1\) 下，最小值为多少？求出最小值点
\((x_1,x_2,x_3)\)。

\subsection{7.}\label{section-5}

线性空间 \(V\) 上的线性变换 \(\mathcal A\) 在基
\((\varepsilon_1,\varepsilon_2,\varepsilon_3,\varepsilon_4)\) 下的矩阵为

\[
A=\begin{pmatrix}
1&0&2&1\\
-1&2&1&3\\
1&2&5&5\\
2&-2&1&-2
\end{pmatrix}.
\]

（1）求 \(\mathcal A\) 的像与核。

（2）在 \(\mathcal A\) 的核中取一组基，把它扩充为 \(V\) 的一组基，并求
\(\mathcal A\) 在这组基下的矩阵。

（3）在 \(\mathcal A\) 的值域中取一组基，把它扩充为 \(V\) 的一组基，并求
\(\mathcal A\) 在这组基下的矩阵。

\subsection{8.}\label{section-6}

若 \(A\) 为 \(n\) 阶实对称矩阵，

\[
W=\{x\in\mathbb R^n\mid x^{\mathsf T}Ax=0\}.
\]

（1）若 \(A\) 为半正定矩阵，证明：\(W\) 是 \(\mathbb R^n\)
的子空间，并求 \(W\) 的维数。

（2）若 \(A\) 为不定矩阵，则 \(W\)
是否满足（1）的结论？若满足，给出证明；若不满足，举出反例。
\end{document}
